\documentstyle[%


righttag%


,11pt%


,amssymb%


,russian%               remove in Eng.


,izvru%                 Set izven% in Eng.


]{amsart}





%       Math definitions





%%%% !!! $Б(K)$ ?? B is used already


\newcommand{\Imm}{\operatorname{Im}}


\newcommand{\const}{\operatorname{const}}


\newcommand{\Ree}{\operatorname{Re}}


\newcommand{\Cc}{\bold C}


\newcommand{\dd}{\mathrm d}


\newcommand{\tts}[1]{}





%\hoffset=-15mm%                  For Eng.


\setcounter{page}{63}


\god{2001}


\nomer{6~(469)} %                        Remove in Eng.


%\nomer{Vol.\,45, No.\,6}%               Set in Eng.


%\str{\pageref{begin}--\pageref{end}}%  Set in Eng.


\udk{517.929}





\author{\it М.А.\,Солдатов, Т.М.\,Митрякова, В.А.\,Ильичев}


% NB:   ^^ remove in Eng.


\title{Отыскание решений линейных дифференциально-разностных уравнений\\


в вырожденном случае}





\date{}


%\pagestyle{myheadings}%         Set in Eng.


\pagestyle{plain} %              Remove in Eng.


\begin{document}


%\thispagestyle{plain}%          Set in Eng.


\maketitle


\label{begin}








{\bf 1.} В статье исследуется вопрос о частных аналитических решениях


линейных неоднородных дифференциально-разностных уравнений с


линейными коэффициентами


\begin{equation}


L[y(x)]\equiv\sum_{i=0}^n \sum_{k=0}^m(a_{ik}x+b_{ik})y^{(i)}(x+h_{k})=F(x),


\end{equation}


где $x$ --- комплексное переменное, $x\in {\bold C}$, а разности


$h_{k}$ --- вещественные числа, причем


\begin{equation}


h_{0}<h_{1}<\dotsb<h_{m}.


\end{equation}


Этот вопрос важен уже потому, что знание частного решения позволяет


свести изучение решений такого уравнения к изучению решений соответствующего


однородного уравнения


\begin{equation}


L[z(x)]=0.


\end{equation}





В [1] установлено, что если $F(x)$ --- целая функция, то уравнение (1) всегда


имеет решение, аналитическое в некоторой правой полуплоскости


$\Ree x>\gamma_{1}$ (в левой полуплоскости $\Ree x<\gamma_{2}$), однако при


этом существенно использовалось условие


\begin{equation}


a_{n0}a_{nm}\ne 0.


\end{equation}


В противоположность ему случай $a_{n0}a_{nm}=0$ условно называем вырожденным.


Вопрос о существовании частных решений


при этом условии остался невыясненным ---


здесь это делается при некоторых специальных


ограничениях на коэффиценты $a_{ik}$, $b_{ik}$


и только в случае, когда $F(x)$ есть целая


функция экспоненциального типа (э.\,т.),


т.\,е. функция из класса $E[1,\infty)$.


(Через $E[1,\sigma]$, $E[1,\infty)$, $E\{1,\sigma\}$, $E\{1,\infty\}$


обозначаем классы целых функций э.\,т.


$\le\sigma$, любого э.\,т., точно порядка 1 и


типа $\sigma$ или $\infty$.)





Для разностного уравнения (1) при $n=0$ изучаемый вопрос полностью


выяснен в ([2], гл.\,3).


При $n\ge1$ получение соответствующих вспомогательных оценок


значительно усложняется.


\medskip





{\bf 2.} Пусть функция $F(x)$ представляется интегралом Лапласа


\begin{equation}


F(x)=\frac{1}{2\pi i}\int_{S}{\phi(t)e^{xt}\dd t},


\end{equation}


где $S$ --- некоторый контур интегрирования (возможно бесконечный) и функция


$\phi(t)$ аналитична вдоль него. Будем искать решения уравнения (1) в форме


такого же интеграла


\begin{equation}


y(x)=\frac{1}{2\pi i}\int_{S}{\gamma(t)e^{xt}\dd t},


\end{equation}


где новую искомую функцию $\gamma(t)$ считаем аналитической вдоль $S$.


Подставляя (6) и (5) в уравнение (1), после очевидных преобразований


(см. [3], [2]) придем к требованию


\begin{equation}


\int_{S}{\Big( T[\gamma(t)]+\phi(t) \Big) e^{xt} \dd t}=


A(t)\gamma(t)e^{xt}\big|_{S},


\end{equation}


где


\begin{gather*}


T[\gamma(t)]=A(t)\gamma'(t)+\big( A'(t)-B(t) \big)\gamma(t),\\


%


A(t)=\sum_{i=0}^{n} {\sum_{k=0}^{m}{a_{ik}t^{i}e^{h_{k}t}} },


\quad B(t)=\sum_{i=0}^{n} {\sum_{k=0}^{m}{b_{ik}t^{i}e^{h_{k}t}} }.


\end{gather*}


$A(t)$ и $B(t)$ называем определяющими функциями (первой и второй) уравнения


(1); $A(t)$ называется обычно характеристической функцией.





Будем обозначать также


\begin{equation}


\gamma_{0}(t)=\frac{1}{A(t)\eta(t)},


\quad \eta(t)=\exp{ \bigg\{ -\int_{t_{0}}^{t}{\frac{B(p)}{A(p)}}\dd p


\bigg\} }


\end{equation}


--- решения дифференциальных уравнений


$$


T[\gamma_0(t)]=0,\quad A(t)\eta'(t)+B(t)\eta(t)=0.


$$


Здесь $l[t_0,t]$ (так обозначаем какой-либо простой путь, соединяющий точки


$t_0$ и $t$) и все рассматриваемые в дальнейшем пути интегрирования не


должны, вообще говоря, проходить через нули функции $A(t)$ --- только эти


нули могут быть особыми точками функций (8).





Как и в [3], [2], устанавливается





\begin{thm}


Пусть $S$ --- конечный замкнутый контур и на нем нет нулей функции $A(t)$.


Тогда для того чтобы функция $(6)$ была решением уравнения $(1)$,


необходимо и


достаточно, чтобы функция $\gamma(t)$ была однозначным на $S$ решением


дифференциального уравнения


\begin{equation}


T[\gamma(t)]=v_1(t)-\phi(t),


\end{equation}


в котором $v_1(t)$ --- какая-либо функция, регулярная на и внутри $S$.


\end{thm}





Общее решение уравнения (9) есть


$$


\gamma(t)=\gamma_0(t) \bigg\{ E+\int_c^t{(v_1(q)-\phi(q))\eta(q)\dd q}


\bigg\},


\quad E=\const.


$$





Пусть $K$ --- выпуклый компакт открытой плоскости {\bf C} (в частности ---


одна точка). Через $Б(K)$ будем обозначать класс целых функций э.\,т.,


сопряженная диаграмма которых лежит в $K$. Через $S(c)$ обозначаем


замкнутый контур с началом и концом в точке $c$, охватывающий компакт $K$


и такой, что в $\overline J(S(c))\backslash K$


нет нулей функции $A(t)$, и считаем, что


$\phi(t)$ в (5) есть функция,


ассоциированная по Борелю с $F(x) \in Б(K)$,


все ее особые точки лежат в $K$.


($J(S) \equiv \,\text{int}\,S$ --- внутренность контура


$S$, $\overline J(S)=J(S)\bigcup S$ и, как обычно, $H(M)$ означает класс


функций, голоморфных на множестве $M$.)





\begin{thm}[{\series{m}\shape{n}\selectfont [2]}]


$1.$ Если $\eta(t)\notin H(K)$, то уравнение $(1)$


имеет решение из класса $Б(K)$


для всякой правой части $F(x) \in Б(K)$.





$2.$ Если $\eta(t)\in H(K)$ , то для существования решения $y(x)\in Б(K)$


необходимо и достаточно условия


\begin{equation}


K_1\equiv\frac{1}{2 \pi i}\int_{S(c)}{\phi(t)\eta(t)\dd t}=0.


\end{equation}


\end{thm}





Эта теорема справедлива и при комплексных $h_k$ и для соответственных


дифференциальных уравнений бесконечного порядка.


Для случая $K=\{t:|t|\le\sigma\}$


другими методами ee


установили Хильб, Перрон, Шеффер, Хапланов (см.~[2], c.\,61).





Пусть $F(x) \in Б(K)$, $\eta(t) \in H(K)$, и возьмем точку $\alpha \in K$ ---


она будет корнем $\eta(t)$ кратности


$s=-\text{res}\,(B(t)/A(t);\alpha)$, $s \ge 0$.


Тогда для уравнения


$$


L[w(x)]=F(x)-\varkappa x^s e^{\alpha x},


\quad \varkappa=K_1/\eta^{(s)}(\alpha),


$$


выполняется условие вида (10), поэтому оно имеет решение $w(x)\in Б(K)$.


При этом заменой $y(x)=w(x)+u(x)$ уравнение (1) приводится к виду


\begin{equation}


L[u(x)]=\varkappa x^s e^{\alpha x}.


\end{equation}





Таким образом, отыскание частных решений уравнения (1) сводится к


отысканию решений уравнения (11). Однако ход доказательства теоремы 2 не дает


эффективного способа нахождения решений, поэтому будем рассматривать


непосредственно исходное уравнение (1) и дадим более конструктивный


метод построения частных решений, причем в обоих случаях: когда $\eta(t)$


регулярна и когда нерегулярна на компакте $K$.





Отметим, что для случая произвольных 


целых функций $F(x)$ (без ограничения на рост)


аналогичная теореме 2 теорема о нормальной разрешимости уравнения (1)


в классе $H({\bold C})$, однако при обязательных условиях (2) и (4),


установлена разными методами в [4] и [1]. При этом в [1] рассмотрены


также решения типа $\Gamma$ (и типа $\Gamma^*$): так называем решения,


регулярные в какой-либо правой полуплоскости $\Ree x>\gamma_1$


(левой полуплоскости $\Ree x<\gamma_2$), но не являющиеся целыми


функциями.


\medskip





{\bf 3.}


Итак, пусть выполнены условия (2) и $F(x)\in Б(K)$. Если $S$ --- конечный


замкнутый контур и решение $\gamma(t)$ уравнения (9) с $v_1(t)\in


H(\overline J(S))$ является функцией, неоднозначной на $S$, то в силу


теоремы 1 интеграл (6) не может определять решение уравнения (1).


Однако, если за $S$ взять бесконечную петлю с ветвями, параллельными


действительной оси (т.\,е. если точку $c\in S=S(c)$ устремить в бесконечность),


то интеграл (6) может дать решение. Пусть $L(c)$ --- путь, который,


начинаясь в точке $c$, удаляется в бесконечность по некоторому лучу


$\{\Imm t=\omega,\ \Ree t\le a \}$ и не проходит через нули


функции $A(t)$, а $L^-(c)$ --- тот же путь, но проходимый в обратном


направлении. Будем рассматривать петли $U(\omega,K)=L^-(c)+S(c)+L(c)$.


Полагаем $x=x_1+ix_2$, $t=t_1+it_2$. Условия на коэффициенты


\begin{equation}


a_{n0}=\dotsb=a_{\nu+1,0}=0, \quad a_{\nu 0}\ne 0;\quad


b_{n0}=\dotsb=b_{\nu+1,0}=0; \quad 0\le\nu\le n,


\end{equation}


выделены в [1] как необходимые условия существования решений типа $\Gamma$


у уравнения (1) с $F(x)\in H(\Cc)$ (в [3]


для уравнения (3) неточно указано условие $a_{n0}\ne 0$, т.\,к. должен


быть частный случай (12) при $\nu =n$).


Решения типа $\Gamma$ являются функциями,


аналитическими во всей плоскости с разрезом по горизонтальному лучу


$U_1=(-\infty+i\Imm\alpha^*,\alpha^*]$, причем точка


$\alpha^*=-b_{\nu 0}/a_{\nu 0}+h_0$ обязательно особая. Покажем, что в


случае (12) при условии $\eta(t)\notin H(\Cc)$ решения типа $\Gamma$


действительно существуют. Их найдем в виде (6), где $S$ --- некоторая


бесконечная петля.





Сначала рассмотрим однородное уравнение (3).


Пусть $\eta(t)\notin H(\alpha_0)$,


причем ограничимся только случаем,


когда $\alpha_0$ --- простой нуль функции $A(t)$. Будем 


полагать $\phi_1(t)=\gamma_0(t)$, когда функция $\gamma_0(t)$


(см. (8)) неоднозначна в окрестности точки $\alpha_0$, или


$\phi_1(t)=\gamma_0(t)\ln(t-\alpha_0)$, когда $\gamma_0(t)\in H(\alpha_0)$.


Утверждаем, что функция


\begin{equation}


Z_1(x)=\frac{1}{2\pi i}\int_{U(0,\{ \alpha_0 \})}{\phi_1(t)e^{xt}\dd t}


\end{equation}


и даст решение типа $\Gamma$ уравнения (3). Для такой функции интеграл в (7)


с $\gamma(t)=\phi_1(t)$ и $\phi(t)\equiv 0$ даст нуль, ибо


$T[\phi_1(t)]\in H(\alpha_0)$. Проверим, что и правая часть обратится


в нуль, и что интеграл (6) сходится в некоторой полуплоскости. Для этого


оценим функцию $\gamma_0(t)$ при вещественных $t\to -\infty$. Учитывая, что


$\frac{B(t)}{A(t)} \to \frac{b_{\nu 0}}{a_{\nu 0}}$, выделим далее из


бесконечно малой $\frac{B(t)}{A(t)} - \frac{b_{\nu 0}}{a_{\nu 0}}$


главную часть вида $\beta/t$, $\beta=\const$. Получим


\begin{equation}


\frac{B(t)}{A(t)}=\frac{b_{\nu 0}}{a_{\nu 0}}+\frac{\beta}{t}+\xi(t),\quad


|\xi(t)|<\frac{B_1}{|t|^2},\quad t\le c_1 <0,


\end{equation}


причем $\xi(t)=B_1(t)/A(t)$, где $B_1(t)$ --- тоже квазиполином вида $B(t)$;


$B_1=\const>0$. При этом учтено, что


\begin{equation}


A(t)=a_{\nu 0}t^{\nu}e^{h_0t}(1+\eta_1(t)), \quad \eta_1(t)\to 0 \quad


\text{при} \quad t\to -\infty.


\end{equation}


Проинтегрировав равенство (14) по пути $l[t_0,t]$, найдем


$$


\int_{t_0}^t\frac{B(p)}{A(p)}\dd p=\frac{b_{\nu 0}}{a_{\nu 0}}t+


\beta\ln t+A_1+\xi_1(t),\quad \xi_1(t)=\int_{t_0}^t\xi(p)\dd p,


$$


$A_1=\const$. К пути $l[t_0,t]$, где считаем $t\le c_1$, добавим двойной луч


$(-\infty,t]$ (т.\,е. луч, проходимый дважды, в разных направлениях). Получим


$$


\xi_1(t)=A_2+\xi_2(t),\quad A_2=\const,


\quad \xi_2(t)=\int_{-\infty}^t\xi(p)\dd p,


$$


$|\xi_2(t)|<B_2/|t|$, $t\le c_1<0$, $B_2=\const>0$. В результате для функции


$\gamma_0(t)$ найдем асимптотику


\begin{equation}


\gamma_0(t)=A_3t^{\beta-\nu}(1+\eta_2(t))


\exp{\bigg\{ \bigg( \frac{b_{\nu 0}}{a_{\nu 0}}-h_0 \bigg) t \bigg\}},


\end{equation}


$\eta_2(t)\to 0$ при $t\to -\infty$, $A_3=\const\ne 0$.





В силу (15) и (16) подстановка в правой части (7) с $S=U(0,\{\alpha_0\})$,


$\gamma(t)=\phi_1(t)$ обращается в нуль, если $\Ree x>\Ree\alpha^*-h_0$,


а интеграл (13) сходится при $\Ree x>\Ree\alpha^*$ и для функции $Z_1(x)$


точка $x=\alpha^*$ является особой, так что $Z_1(x)$ действительно


является решением типа $\Gamma$ уравнения (3). Случай $\nu=n$ изучен в [3].


Отметим, что если $\gamma_0(t)\in H(\alpha_0)$, то интеграл (13) с точностью


до постоянного множителя можно преобразовать в интеграл по простому пути


$l[\alpha_0,-\infty]$ с началом в точке $\alpha_0$ и уходящему в бесконечность


по отрицательной части действительной оси (криволинейный луч), причем от


функции $\phi_1(t)=\gamma_0(t)$.





Вернемся к неоднородному уравнению (1)


с $F(x)\in Б(K)$. Далее 


будем обозначать решение уравнения (9) с $v_1(t)\equiv0$ через $\gamma_1(t)$


\begin{equation}


\gamma_1(t)=-\gamma_0(t)\int_c^t\phi(q)\eta(q)\dd q,


\end{equation}


$T[\gamma_1(t)]=-\phi(t)$. Оценим его при вещественных $t\to -\infty$.


При условии (12) $\forall t\le c_1$ будет $|B(t)/A(t)|\le q_1=\const$,


$q_1\ge |b_{\nu 0}/a_{\nu 0}|$, поэтому


$$


|\gamma_1(t)|\le M_1|A(t)|^{-1}e^{2q_1|t|}(|t|+\chi);


\quad M_1\ \text{и}\ \chi=\const>0.


$$


Тогда правая часть в (7) при $S=U(0,K)$ и $\gamma(t)=


\gamma_1(t)$ обратится в нуль, если $x_1>2q_1$, тождество (7)  будет


выполнено, и функция


\begin{equation}


Y_1(x)=\frac{1}{2\pi i}\int_{U(0,K)}\gamma_1(t)e^{xt}\dd t


\end{equation}


в полуплоскости $x_1>2q_1$ определит решение уравнения (1), причем


интеграл (18) сходится и определяет аналитическую функцию в полуплоскости


$x_1>2q_1+h_0$. В таком случае, как это следует из общих свойств решений


уравнений (1) [1], существует постоянная $C=C_1$ такая, что для


разности $Y_2(x)=Y_1(x)-CZ_1(x)$ точка $x=\alpha^*$ будет правильной


и, следовательно, $Y_2(x)$ определит целое решение уравнения (1).


Тогда при $C\ne C_1$ функции $Y_2(x)$ будут решениями типа~$\Gamma$.


\medskip





{\bf 4.}


Существование решений типа $\Gamma$ уравнений (3) и (1) было установлено


при условии (12) лишь в случае $\eta\notin H(\Cc)$. Покажем, что, независимо


от последнего, решения типа $\Gamma$ всегда существуют при условиях


\begin{gather}


a_{im}=\dotsb=a_{i,s+1}=0 \quad \forall i=0,1,\dotsc,n, \quad


a_{ns}=\dotsb=a_{\lambda+1,s}=0,\notag \\


a_{\lambda s}\ne 0,\quad 0\le\lambda\le n, \quad s<m; \\


b_{nm}=\dotsb=b_{\mu+1,m}=0, \quad b_{\mu m}\ne 0, \quad 0\le\mu\le n.\notag


\end{gather}





Условия (19) при $\mu=n$ были выделены в [5] при построении для уравнения (3)


решений типа $\Gamma$ (если $a_{n0}\ne 0$) и так называемых


элементарных решений второго рода --- они из класса $E\{1,\infty\}$.





Выделяя в два этапа из отношения $B(t)/A(t)$ главные части при


$t_1\to +\infty$, найдем


$$


\frac{B(t)}{A(t)}=


\frac{b_{\mu m}}{a_{\lambda s}}t^{\mu-\lambda}e^{(h_m-h_s)t}+


\frac{q_2}{a_{\lambda s}}t^{\mu-\lambda-1}e^{(h_m-h_s)t}+\alpha_1(t),


$$


где $q_2=\const$, $\alpha_1(t)=O(t^{\mu-\lambda-2})e^{(h_m-h_s)t}$ при


$t_1\to+\infty$. Проинтегрировав это равенство по некоторому пути $l[t_0,t]$,


причем первое слагаемое справа --- по частям два раза, а второе --- один раз,


получим


$$


\int_{t_0}^t\frac{B(p)}{A(p)}\dd p=Bt^{\mu-\lambda}e^{(h_m-h_s)t}+


\alpha_2(t), \quad B=\frac{b_{\mu m}}{(h_m-h_s)a_{\lambda s}}\equiv|B|


e^{i\theta},


$$


причем функция $\alpha_2(t)$, как и ее производная $\alpha'_2(t)$,


на лучах {$T(\tau)=\{ t_2=\tau$, $b(\tau)\le t_1<+\infty \}$,


$\tau$ и $b(\tau)$ --- постоянные, при $t_1\to +\infty$ ведет себя как


$O(t^{\mu-\lambda-1}e^{(h_m-h_s)t})$. Обозначив


$$


u(t_1,t_2)=\Ree \int_{t_0}^t \frac{B(p)}{A(p)} \dd p ,


$$


оценим функцию $\gamma_1(t)$


\begin{equation}


|\gamma_1(t)|\le\frac{1}{|A(t)|}e^{u(t_1,t_2)}\int_{t_0}^t


e^{-u(t_1,t_2)}|\phi(t)|\, |\dd t|.


\end{equation}


Разбив путь $l[t_0,t]$ на три части $l[t_0,\tau]$, $l[\tau,\tau+it_2]$,


$[\tau+it_2,t_1+it_2]$, получим


$$


u(t_1,t_2)=|t|^{\mu-\lambda}e^{(h_m-h_s)t_1}(1+\xi(t_1,t_2))|B|


\cos[\theta+(h_m-h_s)t_2+(\mu-\lambda)\arg t]+\zeta(t_2),


$$


где $\xi(t_1,t_2)\to 0$ при $t_1\to +\infty$, $\forall \tau=t_2=\const$,


$|\zeta(t_2)|\le q_3|t_2|$, $q_3=\const$.





Определим число $\tau$ из условия $\cos[\theta+(h_m-h_s)\tau+


(\mu-\lambda)\arg t]<0$. Так как $\arg t\to 0$ при $t_1\to +\infty$,


$\forall t_2=\const$, то при достаточно большом $b(\tau)$ функция


$(-u(t_1,t_2))$ достигает наибольшее значение в правом конце пути


$l[t_0,t]$, $t\in T(\tau)$. Вынося это значение за знак интеграла в (20),


учитывая ограниченность функции $\phi(t)$, получим


\begin{equation}


|\gamma_1(t)|\le M_1 \frac{|t|+a}{|A(t)|}, \quad M_1=\const, \quad


a=\const>0,\quad t\in T(\tau).


\end{equation}


Для функции же $\gamma_0(t)$ из (8), где полагаем $t_0=c$, на лучах


$T(\tau)$ будет иметь место оценка


\begin{equation}


|\gamma_0(t)|\le Me^{\mu_1 |t|} \exp \{ -d|t|^{\mu-\lambda}


e^{(h_m-h_s)t_1} \}, \quad t_1\ge b(\tau),


\end{equation}


где $\mu_1>0$, $d>0$ --- некоторые постоянные. Оценки (21) и (22)


устанавливаются по той же схеме, как, например, в [1], [5], но технически


несколько сложнее. Учитываем также асимптотику


$$


A(t)=a_{\lambda s}t^{\lambda}e^{h_st}(1+o(1)), \quad t_1\to +\infty.


$$





Пусть $H$ --- петля, состоящая из замкнутого пути $S(c)$ и двойной кривой,


начинающейся в точке $c$, бесконечная часть которой идет по лучу $T(\tau)$.


Тогда интеграл


\begin{equation}


W_1(x)=\frac{1}{2\pi i}\int_H \gamma_1(t)e^{xt}\dd t


\end{equation}


сходится во всяком случае в некоторой полуплоскости $x_1<\rho_0$ и


определяет решение уравнения (1). В силу выполнимости условий (19) всякое


такое решение является целой функцией. Если кроме (19) выполнены и


условия (12), то интеграл


\begin{equation}


Z_2(x)=\frac{1}{2\pi i}\int_{L_1}\gamma_0(t)e^{xt}\dd t,


\end{equation}


где $L_1=(-\infty,c_1]+l[c_1,\tau+ib(\tau)]+[\tau+ib(\tau),


+\infty+ib(\tau)]$, определяет решение типа $\Gamma$ уравнения (3), и поэтому


функции $W_2(x)=W_1(x)+C_2Z_2(x)$ с $C_2=\const\ne 0$ определяют решения типа


$\Gamma$ уравнения (1). (В (24) за путь $L_1$ можно брать и прямую $\Imm t=


\tau$, обходящую нули функции $A(t)$.) При тех же условиях (19) и (12)


функции $Y_3(x)=Y_1(x)-CZ_2(x)$ при некотором значении $C=C_3$ определяют


целые решения уравнения (1), а при $C\ne C_3$ --- решения типа $\Gamma$.





Пусть $\eta(t)\in H(\Cc)$. B силу теоремы 2 решение (23) растет при


$x \to \infty$ быстрее всякой целой функции э.\,т. Можно утверждать, что это


будет функция первого порядка максимального типа (т.\,е. функция из класса


$E\{1,\infty \}$). Для доказательства, заметив, что $\eta(t)\in H(K)$


для любых компактов $K$, разобьем путь $H$ на три части:


$H_1=(+\infty+i\tau,b+i\tau)$ --- верхняя часть петли $H$, $H_2=H_1^-$


--- нижняя часть и $H_3$ --- остальная часть пути $H$; $H_3$ --- это конечная


петля, охватывающая $K$, с началом и концом в точке $t=\tau+ib(\tau)$.


Значение выбранной непрерывной вдоль $H$ ветви функции $\gamma_1(t)$


в точке $t\in H_2$ после обхода по $H$ обозначим $\gamma_1^*(t)$.


В силу однозначности функций $\eta(t)$ и $\gamma_0(t)$, из (17) найдем


$$


\gamma_1^*(t)=\gamma_1(t)-2\pi iK_1\gamma_0(t).


$$


Тогда интеграл (23) перепишется в виде


\begin{equation}


W_1(x)=\frac{1}{2\pi i}\int_{H_3}\gamma_1(t)e^{xt}\dd t-


K_1\int_{H_2}\gamma_0(t)e^{xt}\dd t.


\end{equation}


Интеграл по конечной петле $H_3$ определяет целую функцию из класса


$E[1,\sigma]$, где $\sigma=\max|t|$ $\forall t\in H_3$, а интеграл по лучу


$H_2$ (его обозначим через $W_1^*(x)$)


можем оценить, используя (22), по аналогии


с [5]. В результате получим


$$


|W_1^*(x)|\le M\exp\{ d|x|+qx_1\ln|x_1|\},


$$


где $q=1/(h_m-h_s)$, если $x_1\ge 1$, или $q=0$, если $x_1\le 1$. Отсюда


и из теоремы 2 следует, что  $W_1(x)\in E\{ 1,\infty \}$, если $K_1\ne 0$.





Отметим, что при условии $\eta(t)\in H(K)$ для интеграла (18)


также справедливо


представление вида (25), в силу чего при $K_1=0$ будем иметь


$$


Y_1(x)=W_1(x)=\frac{1}{2\pi i}\int_{S(c)}\gamma_1(t)e^{xt} \dd t


$$


--- это будут целые функции из класса $Б(K)$, что соответствует теореме~2.


\medskip





{\bf 5.}


Резюмируя сказанное в пп.\,3--4, приходим к следующим результатам.





\begin{thm}


В случае $(12)$ при выполнении либо $\eta(t)\notin H(\Cc)$,


либо $(19)$ однородное уравнение $(3)$ имеет решения типа~$\Gamma$.


\end{thm}





\begin{thm}


Пусть $F(x)\in Б(K)$. Тогда


\begin{itemize}


\item[1.] при условии $(12)$ уравнение $(1)$


всегда имеет решение, регулярное в


правой полуплоскости $\Ree x>\Ree \alpha^*$, причем в случае $\eta(t)\notin


H(\Cc)$ существуют как решения типа $\Gamma$, так и целые решения\/$;$


\item[2.] при условии $(19)$ уравнение $(1)$ всегда имеет решение в виде целой


функции, а при дополнительном условии $(12)$ существуют и решения типа


$\Gamma;$


\item[3.] при совместном выполнении условий $\eta(t)\in H(\Cc)$ и $(19)$ в


случае $K_1\ne 0$ целые решения уравнения (1) растут быстрее всякой функции


из класса $E[1,\infty)$ и существуют решения из класса $E\{ 1,\infty\}$.


\end{itemize}


\end{thm}





Отметим, что при условии (19) уравнение (1) не вырождается в обычное


дифференциальное уравнение, и поэтому к последним не относятся


соответствующие положения теоремы 4.


\medskip





{\bf 6.}


В теоремах 3 и 4 речь идет о решениях, аналитических в какой-либо правой


полуплоскости $\Ree x>\gamma_1$. Tакие решения связаны


с условиями на коэффициенты $a_{i0}$. Тогда аналогичные факты можно высказать


и о решениях, аналитических в левой полуплоскости $\Ree x<\gamma_2$ --- и


это будет связано с условиями на коэффициенты $a_{im}$. Впрочем, эти задачи


сводятся одна к другой в результате преобразования уравнения (1) с помощью


замен $x+h_m=-z$, $y(-z)=w(z)$.





Отметим часть таких фактов. Для существования решений типа $\Gamma^*$


уравнения (1) с $F(x)\in H(\Cc)$ необходимы условия


\begin{equation}


|a_{im}|+|b_{im}|=0\ \; \forall i=n,n-1,\dotsc,p+1, \quad a_{pm}\ne 0,\quad


0\le p \le n.


\end{equation}


При этом, если $\eta(t)\notin H(\Cc)$, решения типа $\Gamma^*$ на самом


деле существуют. Такие решения являются функциями, аналитическими во всей


плоскости с разрезом по горизонтальному лучу $U_2=[\beta^*,+\infty+i\Imm


\beta^*)$, причем точка $\beta^*=-b_{pm}/a_{pm}+h_m$ обязательно является


особой.





При выполнении условий


\begin{equation}


\begin{array}{c}


a_{i0}=\dotsb=a_{i,r-1}=0\ \; \forall i=0,1,\dotsc,n, \quad


a_{nr}=\dotsb=a_{q+1,r}=0, \quad r>0,\\


a_{qr}\ne 0\  (0\le q\le n); \quad b_{n0}=\dotsb=b_{t+1,0}=0,


\quad b_{t,0}\ne 0 \ (0\le t \le n)


\end{array}


\end{equation}


уравнение (1) с $F(x)\in E[1,\infty)$ всегда имеет целые решения и они при


$K_1\ne 0$ растут быстрее всякой целой функции э.\,т., а при совместном


выполнении условий (26) и (27) всегда существуют и решения типа $\Gamma^*$.


\medskip





{\bf 7.}


Для уравнения (11) с $s=0$ при условии $A(\alpha)\ne 0$ частные


решения можно построить в значительно более простой форме, именно,


\begin{equation}


y(x)=M\int_{S}\gamma_0(t)e^{xt}\dd t.


\end{equation}


Здесь $\gamma_0(t)$ --- заданная функция из (8), а постоянная $M$ и путь


$S$ подлежат определению из условия


\begin{equation}


\left. MA(t)\gamma_0(t)e^{xt}\right|_S=\varkappa e^{\alpha x},


\end{equation}


которое получается так же, как и (7), с учетом


$T[\gamma_0(t)]\equiv 0$.





С помощью полученных выше оценок функции $\gamma_0(t)$ можно проверить,


что условие (29) будет выполнено, например, в следующих случаях.





а) Пусть выполнено условие (12), и простой путь $S=l(-\infty,\alpha]$


идет из $-\infty$ в точку $\alpha$ по лучу $(-\infty,c_1]$ и далее по


конечному участку $l[c_1,\alpha]$, а число


$M=\varkappa/(A(\alpha)\gamma_0(\alpha))\equiv M_0$.


Тогда интеграл (28) определит решение типа $\Gamma$ уравнения (11).





б) Пусть выполнены условия (19), и простой путь $S=S^*$ идет по лучу


$T(\tau)$ из бесконечности в точку $b(\tau)+i\tau$ и от нее далее в точку


$\alpha$. Тогда при том же значении $M=M_0$ интеграл (28) сходится


$\forall x$ и определяет целое решение (уравнения (11) с $s=0$).





Интересно отметить следующее. Поскольку $A(\alpha)\ne 0$ , то число


$K_1\equiv\varkappa \eta(\alpha)\ne 0$, поэтому в случае $\eta(t)\in H(\Cc)$


интеграл (28) определит решение из класса $E\{ 1,\infty \}$. Так


через свойства решений уравнения (1) (см. теорему 4)


заключаем, что функция, определяемая интегралом вида (28) по пути


$S=S^*$, является целой функцией первого порядка максимального типа.


\medskip





{\bf 8.}


Примеры. 1) Пусть $\Delta u(x)=u(x+1)-u(x)$,


$\Delta^ku=\Delta(\Delta^{k-1}u)$. Для разностного уравнения


$x\Delta^mu(x)=1$ выполняется условие (4) ($n=0$). Оно имеет следующее


решение типа $\Gamma$ (с особыми точками $x=0,-1,-2,\dotsc$):


$u(x)=P_{m-1}(x)\Psi(x)$, где


$P_{m-1}(x)=(x-1)(x-2)\dotsb(x-m+1)/(m+1)!$,


$\Psi(x)=\Gamma'(x)/\Gamma(x)$ --- логарифмическая производная гамма-функции


Эйлера $\Gamma(x)$, и решение типа $\Gamma^*$:


$u^*(x)=P_{m-1}(x)\Psi(m-x)$ (с особыми точками $x=m,m+1,\dotsc$).


Решений в классе целых функций нет.





2) Для разностного уравнения $y(x+1)-xy(x)=1$ выполняются условия вида


(12), (19) ($n=0$). Оно имеет решения


$$


y_1(x)=-e\int_{-\infty}^0\exp(-e^t)e^{xt}\dd t \quad \text{и} \quad


y_2(x)=-e\int_{+\infty}^0\exp(-e^t)e^{xt}\dd t.


$$


Из них первое --- типа $\Gamma$, второе --- целое из класса $E\{ 1,\infty \}$.


Это --- неполные гамма-функции ([2], c.\,103), обычно их записывают в виде


интегралов, получающихся из указанных в результате замены $e^t=\xi$.





3) Возьмем уравнение (3) с $B(t)=(\lambda+1)A'(t)$, $\lambda=0,1,2,\dotsc$


Тогда $\gamma_0(t)=(A(t))^\lambda$. Понятно, что решения типа $\Gamma$


могут существовать и когда $a_{n0}=0$, точнее, когда выполняются условия


(12) с $\nu<n$: достаточно взять квазиполином $A(t)$ с таким свойством.


Например, можно рассмотреть уравнение (3) с $A(t)=t^2e^t+t$, $\lambda=0$;


$A(t)=te^{2t}+t^2e^t+t$, $\lambda=0$ (здесь выполняются оба условия: (12) с


$\nu<n$ и (26) с $p<n$), при этом решением будет $y(x)=1/x$. Случаю


$A(t)=te^t+1$, $\lambda=0$ соответствует уравнение


$(x+1)y'(x+1)+y(x+1)+xy(x)=0$ с решением $y(x)=e^{x\alpha_0}/x$,


где $\alpha_0e^{\alpha_0}+1=0$.





\begin{thebibliography}{9}





\bibitem{1}


Солдатов М.А. {\em О некоторых решениях линейных неоднородных


дифференциально-разностных уравнений} // Изв. вузов. Математика. --


1969. -- \No\,12. -- С.83--92.


\bibitem{2}


Миролюбов А.А., Солдатов М.А. {\em Линейные неоднородные разностные


уравнения}. -- М.: Наука, 1986. -- 128 с.


\bibitem{3}


Леонтьев А.Ф. {\em Дифференциально-разностные уравнения с линейными


коэффициентами} // Тр. Горьковск. пед. ин-та. -- 1951. --


\No\,14. -- С.3--30.


\bibitem{4}


Краплин М.А. {\em О целых и мероморфных решениях одного класса разностных


уравнений} // ДАН СССР. -- 1962. -- Т.\,142. -- \No\,5. -- С.\,1011--1014.


\bibitem{5}


Солдатов М.А. {\em О свойствах решений линейных дифференциально-разностных


уравнений} // Сиб. матем. журн. -- 1967. -- Т.\,8. --


\No\,3. -- С.\,669--679.





\end{thebibliography}





\vskip 2cm





\post{\it Нижегородский государственный}{\it Поступила}


\post{\it университет}{19.07.1999}





\label{end}


\end{document}


